\section{Hardware-Sovereign Serving}
\label{sec:sovereign}

Serving an open-weight LLM is, in practice, often equivalent to depending on a
single vendor's software stack (CUDA/NVIDIA). For a system whose remit is
\emph{national law}, that coupling is itself a sovereignty risk: the guarantee
we make to citizens should not rest on the continued goodwill or export policy
of one foreign supplier. We therefore treat the serving substrate as
replaceable.

\subsection{A thin, replaceable indirection}
Our serving layer is a thin \emph{OpenAI-compatible} indirection. Switching
silicon means changing three environment variables---\texttt{LLM\_BASE\_URL},
\texttt{LLM\_API\_KEY}, \texttt{LLM\_MODEL}---while the pipeline, the verifier,
the service, and the tests are all unchanged. Because the trust guarantee lives
in the verification layer and not in the model, the model becomes a commodity
that can be sourced from whatever accelerator is available and appropriate.

\subsection{What runs where, honestly}
Table~\ref{tab:backends} states, without embellishment, what runs where. We
demonstrate the full pipeline \emph{live} on a Qualcomm Cloud~AI~100 inference
accelerator~\cite{qualcommcloudai100} serving Llama-3.1-8B---dedicated
inference silicon whose value proposition is operations-per-watt rather than
raw training throughput. We target AMD Instinct through a portable inference
runtime~\cite{modularmax2024} that already supports it; European accelerators
are stated as \emph{vision}, not measurement. We report real latency and
tokens/s from the live backend, and we deliberately do \emph{not} print a
fabricated energy figure: we claim only the qualitative, well-founded point
that dedicated inference silicon is engineered for high throughput-per-watt,
and we defer a measured perf/watt study to future work (\S\ref{sec:limits}).

\begin{table}[t]
\centering
\caption{Serving backends and honest status. The pipeline is identical across
rows; only the OpenAI-compatible endpoint changes.}
\label{tab:backends}
\small
\begin{tabular}{@{}lll@{}}
\toprule
Backend & Silicon & Status \\
\midrule
Qualcomm Cloud AI 100 & Qualcomm & \textbf{live, proven} \\
Modular MAX & AMD Instinct & target (runtime supports it) \\
Modular MAX & NVIDIA & backward compatibility \\
Self-hosted (open weights) & NVIDIA (cloud) & fallback \\
VSora / Axelera & EU/FR & vision (no access) \\
\bottomrule
\end{tabular}
\end{table}

\subsection{The OpenAI-compatible contract}
The indirection speaks the OpenAI chat-completions API---\texttt{POST
/v1/chat/completions}, bearer authentication, optional server-sent-event
streaming---now a lingua franca implemented by most serving stacks, including
the portable runtime we target for AMD. Concretely, three environment variables
select the backend, and \emph{nothing else} in the codebase is hardware-aware:

\begin{lstlisting}[style=lr]
# Qualcomm Cloud AI 100 (live)
LLM_BASE_URL = https://.../v1
LLM_API_KEY  = ***
LLM_MODEL    = llama-3.1-8b
# to move to AMD via a portable runtime,
# change these three lines only.
\end{lstlisting}

\subsection{Frugality as a design value}
At national scale, throughput-per-watt is not a vanity metric but a cost and
sustainability constraint. Dedicated inference silicon---engineered for serving
rather than repurposed from training---targets exactly this regime, which is
why we prioritize it. We surface a frugality panel in the public demonstration
(\S\ref{sec:demo}); consistent with our stance on honesty, we present it as an
architectural argument and defer a \emph{measured} perf/watt number until we
can produce one under controlled conditions (\S\ref{sec:limits}).
